% !TeX root = print.tex
% Macros shared by both the web and print versions of the blueprint.
% This file is not built directly. Build src/web.tex or src/print.tex instead.

\usepackage{mathtools}
\usepackage{amssymb}
\usepackage{amsthm}
\usepackage{amsmath}
\usepackage{graphicx}
\usepackage{xcolor}
% alltt: verbatim-like display (preserves spaces/newlines) that still
% honors \, {, } so inline math substitutions for Lean's Unicode syntax
% (\N for the source's `ℕ`, \exists for `∃`, etc. -- see
% ../chapter/retrospective.tex) can sit inside a displayed code block.
% Both xelatex (print) and plastex (web, via its own Packages/alltt.py +
% HTML5/alltt.jinja2 renderer) support it out of the box.
\usepackage{alltt}
% NB: don't \usepackage{enumerate} here. The print build already loads
% enumitem (see ../print.tex), which subsumes enumerate's features;
% loading both has caused \enit@after undefined errors during the print
% build in the docgen-action TeX Live image.
%
% Also: don't \usepackage[normalem]{ulem} here unless we actually use
% \sout / \uline / \uuline / \uwave / etc. in the blueprint prose. plastex
% has no Python implementation of ulem and its LaTeX fallback throws
% during package load ("pop from empty list" on a \write primitive), so
% loading it for nothing pollutes the web build's log.

% Theorem environments. We number everything within section so that labels
% read like 2.3 (= section 2, item 3) rather than running globally.
\theoremstyle{plain}
\newtheorem{theorem}{Theorem}[section]
\newtheorem{lemma}[theorem]{Lemma}
\newtheorem{proposition}[theorem]{Proposition}
\newtheorem{corollary}[theorem]{Corollary}
\theoremstyle{definition}
\newtheorem{definition}[theorem]{Definition}
\newtheorem{remark}[theorem]{Remark}
\newtheorem{example}[theorem]{Example}
% Formalization notes: Lean-encoding asides placed outside every other
% environment. They share the theorem counter (one numbering stream) so a
% \cref reads "Formalization note 2.3"; the \crefname is set after cleveref
% loads (see ../web.tex, ../print.tex).
\theoremstyle{remark}
\newtheorem{fmlnote}[theorem]{Formalization note}
\numberwithin{equation}{section}

\usepackage[hidelinks, colorlinks=false, hypertexnames=false]{hyperref}

% Math shortcuts.
\newcommand{\N}{{\mathbb{N}}}
\newcommand{\Z}{{\mathbb{Z}}}
\newcommand{\Q}{{\mathbb{Q}}}
\newcommand{\R}{{\mathbb{R}}}

% Combinatorial-rigidity-specific notation.
\DeclareMathOperator{\edges}{edges}
\DeclareMathOperator{\vertices}{V}
% Edges of G (default) or another graph (optional first arg) induced on a vertex set.
% Use \edgesIn{S} for E_G[S] or \edgesIn[H]{S} for E_H[S].
\newcommand{\edgesIn}[2][G]{E_{#1}[#2]}
\newcommand{\rk}{\operatorname{rk}}     % rank (rigidity matroid)
\newcommand{\KK}[1]{K_{#1}}             % complete graph on #1
